\documentclass[../../../include/open-logic-section]{subfiles}

\begin{document}

\olfileid[hi]{sth}{choice}{countablechoice}
\olsection{गणनीय चयन}

चयन/सुव्यवस्था के किसी उपयोग के बिना यह सिद्ध करना सरल है कि:

\begin{lem}[$\Zminus$ में]
प्रत्येक परिमित समुच्चय का कोई चयन फलन है।
\end{lem}

\begin{proof}
$a=\{b_1,\ldots,b_n\}$ रखें। सरलता के लिए मान लें कि प्रत्येक
$b_i\neq\emptyset$। अतः ऐसी वस्तुएँ $c_1,\ldots,c_n$ हैं कि
$c_1\in b_1,\ldots,c_n\in b_n$। अब
\olref[z][pairs]{prop:pairsconsequences} से समुच्चय
$\{\langle b_1,c_1\rangle,\ldots,\langle b_n,c_n\rangle\}$ का अस्तित्व है;
और यह $a$ के लिए चयन फलन है।
\end{proof}

किंतु अनंत समुच्चयों पर विचार करते ही विषय अधिक धुँधला हो जाता है। उदाहरण
के लिए ऊपर के इस ``न्यूनतम'' विस्तार पर विचार करें:

\begin{defish}
\emph{गणनीय चयन।} प्रत्येक \emph{गणनीय} समुच्चय का कोई चयन फलन है।
\end{defish}

यह चयन की एक विशेष स्थिति है। पता चलता है कि इस सिद्धांत का उपयोग उसके
उपयोग की स्पष्ट जागरूकता के बिना काफ़ी बार किया गया था। यहाँ दो सुंदर उदाहरण
हैं।\footnote{\citet[\S9.4]{Potter2004} और लूका इंकुर्वाती से।}

\begin{ex}
यह एक स्वाभाविक विचार है: किसी भी समुच्चय $A$ के लिए या तो
$\cardle{\omega}{A}$, अथवा किसी $n\in\omega$ के लिए $\cardeq{A}{n}$। यह
उस सहज विचार को कहने का एक ढंग है कि प्रत्येक समुच्चय या तो परिमित है या
अनंत। कैंटर और अनेक अन्य गणितज्ञों ने यह दावा बिना सिद्ध किए किया। हमने
सावधानी से इसे \olref[cardinals][classing]{generalinfinitycharacter} में
सिद्ध किया। किंतु उस प्रमाण में हम $\ZFC$ में काम कर रहे थे, क्योंकि हमने
माना था कि किसी भी समुच्चय $A$ को सुव्यवस्थित किया जा सकता है, और इसलिए
$\card{A}$ के अस्तित्व की प्रत्याभूति है। अर्थात् हमने स्पष्ट रूप से चयन को
माना था।

वास्तव में \citet{Dedekind1888} ने इस दावे का अपना प्रमाण इस प्रकार दिया:

\begin{thm}[$\Zminus + \text{गणनीय चयन}$ में]
किसी भी $A$ के लिए या तो $\cardle{\omega}{A}$ अथवा किसी $n\in\omega$ के
लिए $\cardeq{A}{n}$।
\end{thm}

\begin{proof}
मान लें कि सभी $n\in\omega$ के लिए $\cardneq{A}{n}$। तब विशेषतः प्रत्येक
$n<\omega$ के लिए कोई उपसमुच्चय $A_n\subseteq A$ है जिसके ठीक
$\cardexpo{2}{n}$ अवयव हैं। इस अनुक्रम $A_0,A_1,A_2,\ldots$ का उपयोग करके
हम प्रत्येक $n$ के लिए परिभाषित करते हैं:
\[
	B_n = A_n \setminus \bigcup_{i < n} A_i.
\]
अब निम्न पर ध्यान दें
\begin{align*}
	\card{\bigcup_{i < n}A_n} 
	&\leq \card{A_0} + \card{A_1} + \ldots + \card{A_{n-1}}\\
	&=1 + 2 + \ldots + 2^{n-1}\\
	& = 2^n - 1\\
	& < 2^n = \card{A_n}
\end{align*}
अतः प्रत्येक $B_n$ का कम-से-कम एक सदस्य $c_n$ है। इसके अतिरिक्त $B_n$
परस्पर असंयुक्त हैं; इसलिए यदि $c_n=c_m$, तो $n=m$। किंतु प्रत्येक
$c_n\in A$। अतः फलन $f(n)=c_n$, $\omega\to A$ एक अंतःक्षेप है।
\end{proof}
\noindent 
डेडेकिंड ने यह नहीं बताया कि उन्होंने गणनीय चयन का उपयोग किया था। किंतु क्या
\emph{आपने} उसका उपयोग देखा? फिर देखें। (सच में: \emph{फिर देखें}।)

प्रमाण ने गणनीय चयन का दो बार उपयोग किया। पहले हमने उससे समुच्चयों का अनुक्रम
$A_0,A_1,A_2,\dots$ प्राप्त किया। फिर प्रत्येक $B_n$ से अपना अवयव $c_n$
चुनने के लिए उसका फिर उपयोग किया। इसके अतिरिक्त चयन का यह उपयोग हटाया नहीं
जा सकता। \citet[p.~138]{Cohen1966} ने सिद्ध किया कि यदि हमारे पास चयन का
कोई रूप न हो, तो परिणाम विफल होता है। अर्थात् $\ZF$ के साथ यह संगत है कि ऐसे
समुच्चय हों जो $\omega$ के साथ \emph{अतुलनीय} हों।
\end{ex}

\begin{ex} 
\citeyear{Cantor1878} में कैंटर ने कहा कि गणनीय समुच्चयों का गणनीय सम्मिलन
गणनीय होता है। उन्होंने प्रमाण नहीं दिया, जिससे शायद संकेत मिलता है कि वे
प्रमाण को स्पष्ट मानते थे। हमने सावधानी से इस परिणाम का अधिक सामान्य रूप
\olref[card-arithmetic][simp]{kappaunionkappasize} में सिद्ध किया। किंतु
हमारे प्रमाण ने स्पष्ट रूप से चयन माना था। और कम सामान्य परिणाम के प्रमाण
को भी गणनीय चयन की आवश्यकता है।

\begin{thm}[$\Zminus + \text{गणनीय चयन}$ में]
यदि प्रत्येक $n\in\omega$ के लिए $A_n$ गणनीय है, तो
$\bigcup_{n<\omega}A_n$ गणनीय है।
\end{thm}

\begin{proof}
व्यापकता खोए बिना मान लें कि प्रत्येक $A_n\neq\emptyset$। अतः प्रत्येक
$n\in\omega$ के लिए कोई !!a{surjection} $f_n\colon\omega\to A_n$ है।
$f\colon\omega\times\omega\to\bigcup_{n<\omega}A_n$ को
$f(m,n)=f_n(m)$ से परिभाषित करें। परिणाम निष्पन्न होता है क्योंकि
$\omega\times\omega$ गणनीय है
(\olref[sfr][siz][zigzag]{natsquaredenumerable}) और $f$ एक
!!a{surjection} है।
\end{proof}
\noindent 
क्या आपने गणनीय चयन का उपयोग देखा? उसका उपयोग फलनों का अनुक्रम
$f_0,f_1,f_2,\dots$ चुनने में होता है।\footnote{चयन का समान उपयोग
\olref[card-arithmetic][simp]{kappaunionkappasize} में हुआ, जब हमने निर्देश
दिया था ``प्रत्येक $\beta\in\cardfont{a}$ के लिए कोई !!a{injection}
$f_\beta$ नियत करें''।} और फिर चयन के किसी सिद्धांत की अनुपस्थिति में परिणाम
विफल होता है। विशेषतः \citet{FefermanLevy1963} ने सिद्ध किया कि $\ZF$ के
साथ यह संगत है कि गणनीय समुच्चयों के किसी गणनीय सम्मिलन की कार्डिनलता
$\beth_1$ हो। किंतु रसेल ने इस बात को कहीं अधिक मज़ेदार ढंग से कहा है:
\begin{quote}
  इसका उदाहरण वह करोड़पति है जो जब भी एक जोड़ी जूते खरीदता था तभी एक जोड़ी
  मोज़े खरीदता था, और किसी अन्य समय नहीं; तथा दोनों खरीदने का उसे इतना शौक
  था कि अंततः उसके पास $\aleph_0$ जोड़ी जूते और $\aleph_0$ जोड़ी मोज़े
  थे\dots\@ जूतों में हम दाएँ और बाएँ का अंतर कर सकते हैं, इसलिए प्रत्येक
  जोड़ी में से एक चुन सकते हैं, अर्थात् सभी दाएँ जूते या सभी बाएँ जूते चुन
  सकते हैं; किंतु मोज़ों के लिए चयन का ऐसा कोई सिद्धांत स्वयं नहीं सूझता,
  और जब तक हम गुणात्मक अभिगृहीत [अर्थात् वस्तुतः चयन] न मानें, हम निश्चित
  नहीं हो सकते कि प्रत्येक जोड़ी में से एक मोज़ा रखने वाला कोई वर्ग है।
  \citep[p.~126]{Russell1919}
\end{quote}
संक्षेप में निम्न सिद्ध करने के लिए चयन के किसी रूप की आवश्यकता है: यदि
आपके पास गणनीयतः अनेक जोड़ी मोज़े हैं, तो आपके पास (केवल) गणनीयतः अनेक
मोज़े हैं। और वास्तव में गणनीय चयन (या किसी तुल्य सिद्धांत) के बिना गणनीय
समुच्चयों का गणनीय सम्मिलन गणनीय होने में विफल हो सकता है।
\end{ex}

निष्कर्ष यह है कि गणनीय चयन का बार-बार उपयोग हुआ, पर उसके उपयोगकर्ताओं को
इसकी अधिक जानकारी नहीं थी। दार्शनिक प्रश्न है: हम गणनीय चयन का
\emph{औचित्य} कैसे दे सकते हैं?

किसी सहज औचित्य का प्रयास अतिसंकार्य का सहारा ले सकता है। मान लें कि हम
पहला चयन $\nicefrac{1}{2}$ मिनट में, दूसरा चयन $\nicefrac{1}{4}$ मिनट में,
\dots, $n$वाँ चयन $\nicefrac{1}{2^n}$ मिनट में, \dots करते हैं। तब एक मिनट
के भीतर हम चयनों का $\omega$-अनुक्रम बना चुके होंगे और एक चयन फलन परिभाषित
कर चुके होंगे।

किंतु ऐसा विचार-प्रयोग वास्तव में हमें क्या बता सकता है? पहली बात, यह
``चुनने'' के विचार को बहुत शाब्दिक रूप से लेने पर निर्भर है। दूसरी बात, यह
गणित को तत्त्वमीमांसात्मक संभावना से बाँधता प्रतीत होता है।

अधिक महत्त्वपूर्ण बात: यह हमें \emph{केवल} गणनीय चयन के बजाय चयन
\emph{सर्वथा} का कोई औचित्य नहीं देगा। क्योंकि यदि हमें \emph{प्रत्येक}
समुच्चय के लिए चयन फलन चाहिए, तो हमें ``मनमानी क्रमसूचक लंबाई का अतिसंकार्य''
कर सकना होगा। स्पष्ट शब्दों में, यह विचार हास्यास्पद है।

\end{document}
